%% ma: L11-B15-0065
%% lop: 11
%% chuong: 5
%% bai: 15
%% dang: 11.15.6
%% loai: TN4
%% muc: VDC
%% dap_an: "C"
%% nguon: "tai-lieu-hoc-tap-mon-toan-lop-11-theo-chuong-trinh-giao-duc-pho-thong-2018.pdf (Lớp Toán Thầy Duy, Phạm Lê Duy, THPT Chu Văn An), Chương 3 Bài 1, Phần 1, Câu 54 (tr.333)"
%% kien_thuc: "Lớp 11 Bài 15 (giới hạn, nguyên lí kẹp); diện tích hình tròn (THCS)"
%% trang_thai: cho-duyet
%% ly_do: "Dạng mới đề xuất (11.15.x), chờ giáo viên duyệt dạng; tài liệu gốc không có lời giải mẫu, đáp án do kho tự giải và kiểm bằng sympy."
%% ghi_chu: "Đã đếm kiểm số: s_n/n^2 -> 3,14159... với n lớn."
\begin{ex}
Với mỗi số nguyên dương $n$, gọi $s_n$ là số cặp số nguyên $(x;y)$ thỏa mãn $x^2+y^2\le n^2$ (nếu $a\ne b$ thì hai cặp số $(a;b)$ và $(b;a)$ khác nhau). Khẳng định nào sau đây là \textbf{đúng}?
\choice{$\lim\dfrac{\sqrt{s_n}}n=\sqrt{2\pi}$.}{$\lim\dfrac{\sqrt{s_n}}n=2$.}{\True $\lim\dfrac{\sqrt{s_n}}n=\sqrt\pi$.}{$\lim\dfrac{\sqrt{s_n}}n=4$.}
\loigiai{\sloppy Mỗi điểm nguyên $(x;y)$ thuộc hình tròn tâm $O$ bán kính $n$ ứng với hình vuông đơn vị tâm tại điểm đó, các hình vuông này không chồng nhau và nằm trong hình tròn bán kính $n+\dfrac{\sqrt2}2$; ngược lại hình tròn bán kính $n-\dfrac{\sqrt2}2$ được phủ bởi các hình vuông ấy. Do đó $\pi\left(n-\dfrac{\sqrt2}2\right)^2\le s_n\le\pi\left(n+\dfrac{\sqrt2}2\right)^2$, suy ra \[\sqrt\pi\left(1-\dfrac{\sqrt2}{2n}\right)\le\dfrac{\sqrt{s_n}}n\le\sqrt\pi\left(1+\dfrac{\sqrt2}{2n}\right).\] Theo nguyên lí kẹp, $\lim\dfrac{\sqrt{s_n}}n=\sqrt\pi$.}
\end{ex}
